\chapter{Síntesis: una empresa, todos los lentes}\label{ch:capstone}

% Alcance: la empresa de referencia B2B SaaS recorrida a través de cada Parte del libro en
%   secuencia, trazando cómo cada disciplina alimenta la siguiente a lo largo de la columna
%   vertebral de creación de valor. Sin nuevas etiquetas eq:. Solo referencias cruzadas.
%   Una figura: value-spine. Estructura: párrafo introductorio → tabla de la hoja de datos →
%   9 secciones (una por Parte + los añadidos de 2026).

Cada \emph{framework} de este libro se introdujo de forma independiente. Aquí convergen. La misma
empresa B2B SaaS de \money{4000000} de ARR que apareció en cada ejemplo trabajado se recorre
de principio a fin a través de las 9 Partes y los 32 capítulos: desde la primera prueba de
creación de valor y el sustrato contable que fundamenta cada número, pasando por la economía
del cliente y la adquisición paga, entrando en la posición competitiva y la durabilidad
estratégica, luego el modelo financiero que valora el conjunto, y finalmente las disciplinas de
ejecución que determinan si el plan sobrevive al contacto con un trimestre. El propósito no es
la repetición --- cada sección supone que su capítulo ha sido leído --- sino la integración:
las conexiones que los capítulos por separado solo podían prometer son ahora explícitas. Las
siete Partes operativas (Fundamentos, Cliente, Adquisición, Estrategia, Finanzas, Ejecución y
el Fundador) van seguidas de una Parte de Referencia y esta Síntesis; los once capítulos
añadidos en la edición de 2026 --- que cubren estados financieros, experimentos, gestión de
producto, canales y salida al mercado (\emph{GTM}), estrategia de datos e IA, el modelo
operativo, la mecánica de la captación de fondos y el propio oficio del fundador en
persuasión, narrativa, psicología y liderazgo --- están integrados en cada sección a
continuación.

\medskip
\noindent\textbf{Hoja de datos de la empresa de referencia.}

\medskip
\noindent\begin{tabularx}{\linewidth}{@{}l >{\raggedright\arraybackslash}X@{}}
  \toprule
  Parámetro & Valor \\
  \midrule
  Precio del plan (ARPU)                 & \money{50}/mes (\money{600}/año) \\
  LTV (perpetuidad creciente)            & \money{3150} \\
  Margen de contribución                 & \qty{42}{\percent} → \money{21}/mes \\
  CAC objetivo (operativo)               & \money{600} → LTV:CAC 5,25× \\
  Crecimiento en estado estacionario $g$ & \qty{3}{\percent}/año \\
  Presupuesto de anuncios mensual → nuevos clientes & \money{120000} → 200/mes \\
  WACC / tasa de descuento               & \qty{11}{\percent}/año \\
  \emph{Funnel}                          & 10.000 sesiones → 800 pruebas → 200 pagos \\
  NOPAT (aprox.)                         & \money{600000} \\
  ARR / MRR                              & \money{4000000} / \money{340000} \\
  Capital invertido                      & \money{3000000} \\
  ROIC                                   & \qty{20}{\percent} \\
  Beta ($\beta$)                         & 1,4 → costo de \emph{equity} \qty{11}{\percent} \\
  Margen bruto                           & \qty{70}{\percent} \\
  Margen EBITDA año 3                    & \qty{+7}{\percent} (puntuación Rule-of-40: 37) \\
  Valoración pre-money Serie A (mercado) & \money{24000000} (\(6\times\) ARR × ajuste Rule-of-40) \\
  Valoración pre-money Serie A (piso DCF) & \money{15000000} (escenario base, \cref{eq:scenario-ev}) \\
  \emph{Runway} al \emph{burn} actual   & 8 meses (\money{2500000} en caja ÷ \money{310000} \emph{burn} neto) \\
  SAFE pendiente                         & \money{500000} a cap de \money{5000000}, descuento \qty{20}{\percent} \\
  \bottomrule
\end{tabularx}
\medskip

% complexity: 5 -> figura value-spine

\begin{figure}[htbp]
  \centering
  \includegraphics[width=0.96\linewidth]{value-spine}
  \caption{La columna vertebral de creación de valor. Cada etapa del libro alimenta la
    siguiente: los Fundamentos establecen la prueba ROIC\,>\,WACC; la economía del Cliente
    fija el CLV que limita el gasto en adquisición; la Adquisición valida la eficiencia del
    canal mediante ROAS e iROAS; la Estrategia establece la tasa de crecimiento $g$ impulsada
    por el \emph{moat} que se usa en el valor terminal; las Finanzas descuentan los flujos de
    caja resultantes al WACC; las disciplinas de Ejecución (Rule of 40, \emph{quick ratio})
    confirman que el plan se está realizando; las habilidades del Fundador (influencia,
    narrativa, liderazgo) convierten el plan en clientes cerrados, una ronda de financiación y
    un equipo dirigido; el capítulo de Referencia cierra el bucle de medición.}
  \label{fig:value-spine}
\end{figure}

%% ─────────────────────────────────────────────────────────────
\section{Fundamento: ¿crea valor?}\label{sec:capstone-value}
%% ─────────────────────────────────────────────────────────────

La prueba maestra es \cref{eq:economic-profit}: el beneficio económico es igual al NOPAT menos
el cargo por capital. El NOPAT de \money{600000}, el capital invertido de \money{3000000} y el
ROIC del \qty{20}{\percent} no son supuestos --- se derivan del modelo de tres estados financieros
de \cref{ch:financial-statements}. \Cref{sec:three-statements} y \cref{sec:statement-analysis}
rastrean cada dólar desde el estado de resultados hasta el puente de NOPAT de
\cref{eq:nopat-bridge}: EBIT \(\times (1-30\%) = \money{860000}\times0.70 \approx \money{600000}\).
Los operadores que no pueden rastrear esos números hasta los estados contables trabajan desde la
fe. La disciplina de inferencia de \cref{ch:experiments-and-inference} subyace a cada afirmación
causal del libro: la maquinaria estadística de \cref{eq:z-test} y \cref{eq:sample-size} determina
si una mejora observada de ROIC es real o ruido muestral, y el \emph{framework} OLS de
\cref{eq:ols} controla las covariables que de otro modo confundirían la atribución.

Con NOPAT \(\approx\money{600000}\) y capital invertido de \money{3000000},
\cref{eq:roic} arroja \(\mathrm{ROIC}=20\%\). El WACC, derivado en \cref{ch:risk-and-capital-structure}
mediante \cref{eq:wacc} y \cref{eq:capm} con \(\beta=1.4\), es \qty{11}{\percent}:
\[
  \mathrm{EP}
    = (0.20 - 0.11)\times\money{3000000}
    = \qty{9}{\percent}\times\money{3000000}
    = \money{270000}.
\]
El diferencial es positivo y sustancial. Una brecha ROIC--WACC del \qty{9}{\percent} no es un
error de redondeo; es evidencia de que el negocio gana \money{270000} anuales por encima del
costo de oportunidad de los inversores. Cada dólar adicional de capital invertido desplegado
al ROIC actual genera \money{0.09} de beneficio económico anual --- una afirmación de
capitalización, no una ganancia única.

La implicación de \cref{eq:economic-profit} opera en ambas direcciones. Si el WACC sube ---
por ejemplo, porque el aumento de las tasas empuja \cref{eq:capm} del \qty{11}{\percent} al
\qty{13}{\percent} --- el diferencial se comprime del \qty{9}{\percent} al \qty{7}{\percent} y el
beneficio económico cae a \money{210000} sin ningún deterioro operativo. A la inversa, reducir el
capital invertido en \money{500000} (mediante una gestión más estricta del capital de trabajo)
manteniendo el NOPAT constante eleva el ROIC al \qty{24}{\percent} y el beneficio económico a
\money{325000}. El denominador del capital es tan importante como el numerador.

Este veredicto de \cref{ch:value-creation} enmarca todo análisis posterior. Los capítulos que
siguen traducen esta única pregunta a sus propios términos: CLV:CAC es su expresión por cliente
(\cref{ch:customer-economics}); el valor de empresa en DCF es su valor presente acumulado
(\cref{ch:valuation}); el \emph{moat} es el mecanismo que lo sostiene
(\cref{ch:advantage-and-disruption}). \textcite{brealey2020} y \textcite{damodaran2012}
desarrollan la arquitectura de finanzas corporativas detrás de \cref{eq:economic-profit,eq:roic,eq:wacc}.

%% ─────────────────────────────────────────────────────────────
\section{Comprensión del cliente: ¿quién y cuánto?}\label{sec:capstone-customer}
%% ─────────────────────────────────────────────────────────────

La creación de valor comienza con la elección de a quién servir. \cref{eq:segment-score} aplicado en
\cref{ch:segmentation-targeting-positioning} puntúa al mercado mediano (50--500 puestos) en 7,70
frente a 6,20 de las pymes, considerando tamaño de mercado, intensidad competitiva, potencial
de LTV y encaje con el ciclo de ventas. El criterio decisivo es el potencial de LTV: el ARR del
mercado mediano es aproximadamente \(10\times\) el nivel de las pymes, lo que se propaga en
cascada a cada número posterior.

Con el segmento objetivo fijado, el margen de precio está acotado por el valor económico para
el cliente (\cref{eq:evc} en \cref{ch:pricing}): la contribución del producto al resultado del
cliente establece el techo de precio teórico por encima del cual la disposición a pagar colapsa.
Nuestro precio de \money{50}/mes se sitúa muy por debajo de ese techo para los compradores de
mercado mediano --- la justificación principal del movimiento de desarrollo de producto hacia un
\emph{tier} empresarial de \money{299}/mes explorado en \cref{ch:growth}.

El caso financiero por cliente es \cref{eq:clv}. Con un margen de contribución del \qty{42}{\percent}
(\money{252}/año), WACC del \qty{11}{\percent} y crecimiento del margen en estado estacionario del
\qty{3}{\percent}:
\[
  \mathrm{CLV} \;=\; \frac{\money{252}}{0.11 - 0.03} \;=\; \money{3150}.
\]
Este es el techo del gasto racional en adquisición. \cref{eq:ltv-cac} lo convierte en la
restricción operativa: con un CAC objetivo de \money{600}, LTV:CAC\,=\,5,25×, muy por encima del
piso de 3× que sirve de proxy para la prueba ROIC\,>\,WACC a nivel de cliente. La brecha entre
5,25× y el piso de 3× representa el margen que absorbe la incertidumbre del modelo en el CLV ---
el equivalente práctico del margen de seguridad en \cref{ch:corporate-finance-core}.

El enfoque de marca de \cref{ch:brand-growth} añade una observación del lado de la demanda: la
notoriedad de marca reduce el esfuerzo mental que los compradores dedican en los puntos de
entrada a la categoría, reduciendo el CAC efectivo en el margen. Una empresa que gana el CLV de
\money{3150} pero gasta solo \money{400} en CAC porque los compradores llegan ya predispuestos
tiene una posición estructuralmente superior en la economía de unidad --- la misma prueba ROIC,
ejecutada con menor costo de capital.

\Cref{ch:product-management} establece \emph{qué} construir antes de que el precio fije cuánto
cobrar. El \emph{tier} empresarial --- la expansión de producto de \money{299}/mes referenciada en
\cref{ch:pricing} --- no surgió de la intuición. Superó la puerta de priorización RICE de
\cref{eq:rice}, que lo puntuó por encima de las alternativas al alcanzar a más usuarios con mayor
confianza por unidad de esfuerzo de ingeniería. El modelo de lanzamiento Bass (\cref{eq:bass})
pronostica luego con qué rapidez el segmento empresarial absorberá el nuevo \emph{tier}: con un
techo alcanzable de 500 cuentas de mercado mediano y un coeficiente de imitación de \(q=0.3\),
el volumen de nuevos adoptantes alcanza su punto máximo alrededor del mes 10, antes del cual la
captación debe concentrarse para sembrar el término social. El encaje producto-mercado en sí
se mide mediante \cref{eq:retention}: la curva de supervivencia asintótica que confirma el
encaje genuino es la misma curva de retención por cohorte que impulsa el CLV.

En conjunto, \cref{ch:segmentation-targeting-positioning}, \cref{ch:brand-growth},
\cref{ch:pricing}, \cref{ch:product-management} y \cref{ch:customer-economics} responden la
mitad de la ecuación de valor correspondiente al cliente: quién genera \money{3150} de valor,
qué construir para sostenerlo, y a qué costo de adquisición.
\textcite{gupta2004customer} es el tratamiento más riguroso del CLV como impulsor del valor de
la empresa.

%% ─────────────────────────────────────────────────────────────
\section{Adquisición: comprar demanda con rentabilidad}\label{sec:capstone-acquisition}
%% ─────────────────────────────────────────────────────────────

La economía del cliente establece cuánto vale una conversión; los canales de adquisición
determinan cuánto cuesta. \cref{eq:funnel-chain} en \cref{ch:digital-funnel-and-crm} traza la
cascada de conversión: 10.000 sesiones producen 800 pruebas gratuitas a una tasa de prueba del
\qty{8}{\percent}, y 200 de esas pruebas se convierten en clientes de pago al \qty{25}{\percent}.
La tasa de conversión global es del \qty{2}{\percent}. La restricción vinculante en el estado
estacionario es la puerta de conversión de prueba a pago, no la parte superior del \emph{funnel}
--- un hallazgo de la teoría de restricciones que \cref{ch:operations} generaliza: elevar el paso
de sesiones no vinculante del \qty{8}{\percent} al \qty{10}{\percent} añade solo 50 clientes
adicionales al mes; elevar la tasa de conversión de prueba a pago del \qty{25}{\percent} al
\qty{35}{\percent} añade 80.

En el lado pago, \cref{eq:roas-cpa} de \cref{ch:paid-acquisition-platforms} rige la prueba de
eficiencia. Un CPA objetivo de \money{600} con \money{120000}/mes entrega 200 conversiones ---
precisamente el volumen en la fase de aprendizaje que los algoritmos de distribución de las
plataformas necesitan para estabilizarse. Fijar el objetivo en \money{300} para «mejorar la
eficiencia» reduce a la mitad el volumen de eventos, atrapa el conjunto de anuncios en Aprendizaje
Limitado y hace que el CPA realizado suba en lugar de bajar: la plataforma no puede estimar
\(\hat{p}\) con precisión por debajo de aproximadamente 50 eventos de conversión por semana.

La pregunta más difícil es si el CPA de \money{600} refleja clientes verdaderamente incrementales
o incluye demanda orgánica que habría convertido de todas formas. \cref{eq:iroas} de
\cref{ch:attribution-and-measurement} lo responde mediante experimentos de geo-retención o
\emph{ghost-bidding}: el iROAS aísla el incremento atribuible a los medios pagos, eliminando la
inflación del último clic. Un negocio que opera con un iROAS sustancialmente inferior a su
supuesto de ROAS objetivo está subvencionando la demanda orgánica con presupuesto pago y
sobreestimando la eficiencia del canal.

\Cref{ch:channels-gtm-and-sales} se sitúa aguas arriba del propio \emph{funnel}: el movimiento
GTM elegido en \cref{sec:gtm-motion} determina qué \emph{funnel} construye la empresa. Un
movimiento de crecimiento liderado por el producto (PLG) llena la parte superior del \emph{funnel}
con activaciones de prueba de autoservicio; un movimiento liderado por ventas lo llena con
transferencias de MQL. Para esta empresa, un movimiento híbrido producto-liderado-ventas es
óptimo --- PLG para pymes, asistido por SDR para PQL de mercado mediano. El dimensionamiento de
la fuerza de ventas fluye a través de \cref{eq:sales-productivity}: para añadir \money{2000000}
de ARR con cuotas de \money{500000} al \qty{80}{\percent} de cumplimiento, se requieren cinco
ejecutivos de cuenta, con un CAC asistido por ventas mixto de aproximadamente \money{180} ---
aún dentro del piso de LTV:CAC de \cref{eq:ltv-cac}. Ese número de personas y el costo de
compensación se incorporan directamente en el modelo operativo de \cref{ch:operating-model}
como partidas en la proyección de gasto en S\&M.

Los tres capítulos de la Parte~III (\cref{ch:digital-funnel-and-crm}, \cref{ch:paid-acquisition-platforms},
\cref{ch:attribution-and-measurement}) forman un bucle cerrado: el \emph{funnel} cuantifica
lo que se convierte; la mecánica de la plataforma determina cuánto cuesta; la atribución
confirma si el gasto fue causal. \Cref{ch:channels-gtm-and-sales} extiende el bucle una etapa
antes, a la elección estructural de \emph{cómo} llegar al mercado antes de comprometer
cualquier dólar de gasto en adquisición.

%% ─────────────────────────────────────────────────────────────
\section{Estrategia: posicionar y defender}\label{sec:capstone-strategy}
%% ─────────────────────────────────────────────────────────────

Una buena economía de unidad en un mercado estructuralmente débil es temporal.
\cref{eq:hhi} en \cref{ch:competitive-analysis} sitúa el mercado B2B SaaS en aproximadamente
\num{2260} (puntos porcentuales al cuadrado), justo por debajo del umbral de alta concentración
--- moderadamente concentrado, con rivalidad significativa y bajas barreras de entrada. En el
lenguaje de las Cinco Fuerzas de Porter (\cref{ch:competitive-analysis}): el poder de los
compradores es alto (los costos de cambio son bajos sin integraciones), la amenaza de entrada
es alta (los costos de infraestructura han caído drásticamente) y la rivalidad es intensa en un
punto de precio de \money{50}/mes donde la paridad de características es habitual. El
diagnóstico estructural es claro: la empresa no puede ganar por precio ni por completitud de
características a su escala actual. Debe construir barreras estructurales.

La taxonomía de \emph{moat} de \cref{ch:growth} identifica los costos de cambio como la palanca
principal. Cada integración API añadida al flujo de trabajo de un cliente aumenta el costo de
migración al salir en un estimado de \money{3000}--\money{6000} --- varios años de valor de
suscripción al precio actual. Un cliente profundamente integrado se enfrenta a una decisión de
deserción genuinamente irracional incluso si llega un producto nominalmente superior. El
\emph{moat} secundario es el efecto de red del mercado de integraciones: \cref{eq:metcalfe}
muestra que con \(n=100\) integraciones de proveedores, existen \(n(n-1)/2 = 4{.}950\) pares
de interacción únicos, frente a solo 1.225 con \(n=50\). Cada nueva integración hace más valiosa
a cada integración existente, lo que atrae a más proveedores, lo que atrae a más clientes ---
la capitalización por el lado de la demanda que está ausente en una ventaja de calidad de
producto pura.

La prueba VRIN de \cref{ch:advantage-and-disruption} aplica: los datos de uso propietarios
acumulados en el mercado de integraciones son \emph{valiosos} (habilitan características de
comparación), \emph{raros} (los competidores no pueden observar el mismo conjunto de datos
entre clientes), \emph{inimitables} (los datos solo se acumulan operando a escala con el tiempo)
y \emph{no sustituibles} (ningún conjunto de datos público replica lo que crea la red).
Un activo de datos que satisface VRIN sostiene el diferencial ROIC\,>\,WACC de \cref{eq:roic}
más allá del horizonte en que cualquier ventaja de característica individual se habría erosionado.

\Cref{ch:data-and-ai-strategy} operacionaliza la prueba VRIN de \cref{sec:rbv-vrin} para la
capa de datos específicamente. La telemetría de uso multicliente propietaria (\cref{sec:data-asset})
supera las cuatro pruebas VRIN: la característica de comparación que habilita es valiosa, el
corpus es raro e inimitable, y ningún conjunto de datos público lo sustituye. La decisión de
construir versus comprar en IA (\cref{sec:ai-build-buy}) aplica \cref{eq:npv} para determinar
que el ajuste fino (\emph{fine-tuning}) sobre este corpus solo tiene NPV positivo si la retención
mejora en aproximadamente \money{50} por cliente, no solo un incremento marginal de calidad. La
gobernanza de IA (\cref{sec:ai-governance}) enmarca el riesgo del modelo como un problema de
agencia y concentra la revisión humana en el bucle donde el producto de la probabilidad de error
y el CLV es mayor. La materialidad ESG (\cref{sec:esg-data}) se trata como un problema de
medición con límites de doble materialidad, donde el mecanismo financiero principal a
\money{4000000} de ARR es la reducción del WACC mediante la reducción de la fricción en la
captación de fondos, no el NOPAT directamente.

La producción financiera de la sección de estrategia se incorpora directamente en la valoración:
el \emph{moat} justifica una tasa de crecimiento terminal $g$ distinta de cero en
\cref{eq:terminal-value}. Sin un \emph{moat}, la competencia comprime $g$ hacia cero; con un
efecto de red de integración y un volante de datos con VRIN positivo, $g$ puede mantenerse en
\qty{3}{\percent} o más durante el período terminal. \textcite{porter1985} fundamenta el
\emph{framework} de análisis competitivo; \textcite{dranove7theconomics} conecta el HHI y la
fortaleza del \emph{moat} con los márgenes sostenibles.

%% ─────────────────────────────────────────────────────────────
\section{Finanzas: el valor cuantificado}\label{sec:capstone-finance}
%% ─────────────────────────────────────────────────────────────

\cref{ch:risk-and-capital-structure} fundamenta la tasa de descuento. \cref{eq:capm} con
\(\beta=1.4\), una tasa libre de riesgo del \qty{4}{\percent} y una prima de riesgo de
\emph{equity} del \qty{5}{\percent} arroja un costo de \emph{equity} del \qty{11}{\percent}.
Como la empresa está financiada con \emph{equity}, \cref{eq:wacc} se reduce a ese mismo
\qty{11}{\percent}. Este único número se propaga por cada cálculo de valor presente del libro:
es el denominador de \cref{eq:clv}, la tasa de descuento en \cref{eq:dcf} y la referencia de
tasa de crecimiento en \cref{eq:terminal-value}.

\cref{eq:dcf} en \cref{ch:valuation} descuenta cinco años explícitos de pronóstico de flujo de
caja libre. El FCF del Año 1 es \money{600000}; la trayectoria crece al \qty{30}{\percent} en
los años 1--3 y luego desacelera al \qty{10}{\percent} en los años 4--5, produciendo una suma de
valores presentes de los años explícitos de aproximadamente \money{3380000}. El valor terminal
es el componente mayor:

\begin{example}[Cálculo del valor terminal]
\cref{eq:terminal-value} con un FCF normalizado del año 5 de \money{1263748} (\money{1226940}
aumentado un período al \qty{3}{\percent}), descontado a \(\mathrm{WACC}-g = 0.11-0.03 = 0.08\):
\[
  \mathrm{TV} \;=\; \frac{\money{1263748}}{0.08} \;\approx\; \money{15800000}.
\]
Añadir el valor presente del valor terminal (\(\approx\money{9400000}\) después de descontar
5 años al \qty{11}{\percent}) a la suma de los años explícitos da un valor de empresa
\(\approx\money{12800000}\). El valor terminal representa el \qty{74}{\percent} del total ---
es una apuesta por el estado estacionario, no por los próximos cinco años.
\end{example}

El análisis de escenarios mediante \cref{eq:scenario-ev} convierte esta estimación puntual en
una distribución ponderada por probabilidades. Tres escenarios --- pesimista (\(p=0.25\), EV
\money{8000000}), base (\(p=0.50\), EV \money{15000000}), optimista (\(p=0.25\), EV
\money{28000000}) --- arrojan:
\[
  EV^* \;=\; 0.25\times\money{8000000}
           + 0.50\times\money{15000000}
           + 0.25\times\money{28000000}
         \;=\; \money{16500000}.
\]
La brecha de \money{3700000} entre el DCF de un solo punto (\money{12800000}) y el EV esperado
(\money{16500000}) refleja el peso del escenario optimista: su diferencial al alza supera el del
pesimista a la baja, de modo que ponderar por probabilidad eleva la expectativa por encima de la
estimación conservadora. Capturar el escenario optimista --- efectos de red de plataforma ---
es tan valioso como gestionar el riesgo del escenario pesimista de erosión del \emph{moat} por
regulación de portabilidad de datos.

\Cref{ch:operating-model} construye los flujos de caja libres explícitos del DCF desde cero en
lugar de tomarlos como dados. \Cref{eq:cohort-revenue} apila la mecánica de supervivencia de
cada cohorte mensual en un pronóstico de ARR --- el propio ARR de \money{4000000} es la
producción de 38 meses de \num{200} nuevos clientes al mes con un \emph{churn} mensual del
\qty{0.65}{\percent}. El modelo EBITDA de tres impulsores de \cref{eq:op-model} convierte luego
el ARR y las tasas de costo en el FCF que alimenta \cref{eq:dcf}. Para el Año~3, con S\&M
cayendo del \qty{50}{\percent} al \qty{35}{\percent} del ARR a medida que el canal madura, el
margen EBITDA alcanza el \qty{+7}{\percent} y la puntuación Rule-of-40 sube a 37 ---
acercándose al umbral de referencia. El análisis de sensibilidad mediante \cref{eq:sensitivity}
muestra que un aumento de \money{100} en el CAC reduce el EBITDA anual en \money{240000};
un aumento del \emph{churn} del \qty{0.35}{\percent} reduce el techo de ingresos en estado
estacionario en \money{1400000} --- un orden de magnitud mayor que un año de gasto excesivo en
CAC. El modelo operativo convierte estas sensibilidades en las trayectorias de FCF pesimista,
base y optimista que \cref{eq:scenario-ev} pondera por probabilidad.

\Cref{ch:fundraising-and-dilution} convierte el valor de empresa DCF de \money{12800000} en
una valoración pre-money de la Serie~A. Los inversores no usan el DCF de forma aislada; aplican
\cref{eq:arr-multiple} para anclar la pre-money a los comparables de mercado. Con ARR de
\money{4000000} y una puntuación Rule-of-40 de 15 (Año~1), el múltiplo de mercado se descuenta
de \(8\times\) a \(6\times\), dando una pre-money implícita de mercado de \money{24000000} ---
cómodamente entre los escenarios base (\money{15000000}) y optimista (\money{28000000}) de
\cref{eq:scenario-ev}. La estimación puntual DCF de \money{12800000} sirve así como piso
conservador; el EV ponderado por probabilidad de \money{16500000} y el múltiplo de mercado
\(6\times\) se sitúan ambos por encima de él. El SAFE pendiente ---
\money{500000} a un cap de \money{5000000}, descuento del \qty{20}{\percent}
(\cref{eq:safe-conversion}) --- convierte al cap, dando al inversor SAFE el \qty{10}{\percent}
pre-Serie~A, antes de que se admita el nuevo dinero institucional. Un \emph{runway} de 8 meses
frente a un proceso de 4--6 meses significa que el disparador de captación de fondos ya se ha
cruzado: \cref{eq:runway} con \money{2500000} en caja y \money{310000} de \emph{burn} neto no
deja margen para el retraso.

\textcite{damodaran2012} ofrece el tratamiento más sistemático de la sensibilidad del valor
terminal; \textcite{brealey2020} cubre el umbral de rentabilidad por escenarios como complemento
al DCF de estimación puntual.

%% ─────────────────────────────────────────────────────────────
\section{Ejecución: restricciones e incentivos}\label{sec:capstone-execution}
%% ─────────────────────────────────────────────────────────────

El valor financiero es una afirmación de valor presente sobre la ejecución futura. El enfoque
operativo pregunta si existen los procesos para materializarlo. \cref{eq:littles-law} en
\cref{ch:operations} establece \(L = \lambda W\): con 40 operaciones abiertas en el
\emph{pipeline} de mercado mediano y 10 cerradas al mes, el ciclo de ventas promedio es
exactamente 4 meses. Comprimirlo a 2 meses requiere o bien reducir a la mitad el
\emph{pipeline} (eliminando operaciones no calificadas) o duplicar la tasa de cierre ---
pero no añadir personal SDR para inflar \(L\) mientras \(\lambda\) es fija, lo que
alargaría \(W\) de forma mecánica.

En el lado de la eficiencia del crecimiento, \cref{ch:business-models} proporciona dos pruebas
en tiempo real. En un mes dado, la empresa SaaS añade \money{30000} en nuevo MRR y
\money{20000} en expansión; pierde \money{5000} por contracción y \money{22000} por
\emph{churn}. \cref{eq:quick-ratio} da:
\[
  \text{Quick Ratio} \;=\; \frac{\money{30000}+\money{20000}}{\money{5000}+\money{22000}}
    \;=\; \frac{50000}{27000} \;\approx\; 1.85.
\]
El crecimiento es positivo pero el denominador (\money{27000} perdidos) casi compensa el flujo
de expansión (\money{20000}), revelando que la base existente está cerca de auto-sabotearse. El
principal cuello de botella es el \emph{churn}, no la adquisición de nuevos clientes.

\cref{eq:rule40} confirma el perfil crecimiento-rentabilidad: con un crecimiento de ARR del
\qty{30}{\percent} y un margen EBITDA de \(\qty{-15}{\percent}\), el Rule of 40 puntúa
\(30 + (-15) = 15\) --- por debajo del umbral del \qty{40}{\percent} pero coherente con una
fase de inversión eficiente en capital. \cref{eq:burn-multiple} con \money{1500000} de caja
neta consumida frente a \money{1200000} de ARR neto nuevo da un múltiplo de \emph{burn} de
\(1.25\times\), que es sólido: el negocio genera más de un dólar de ARR nuevo por cada dólar
y cuarto consumido, muy por debajo del umbral de precaución de 2.

El riesgo organizativo es la desalineación de incentivos que \cref{ch:organizational-behavior}
denomina por el nombre de Kerr: recompensar lo que es medible (MRR bruto nuevo, volumen del
\emph{funnel}) en lugar de lo que importa (beneficio económico neto, ROIC). Un equipo de ventas
compensado por operaciones cerradas no tiene incentivos para evitar a los clientes mal encajados
con alta probabilidad de \emph{churn}; un equipo de producto recompensado por la velocidad de
características no tiene incentivos para construir la profundidad de integración que reduce el
\emph{churn}. La gobernanza del consejo (\cref{sec:board-governance}) proporciona el respaldo
estructural: un consejo constituido con miembros independientes, un comité de auditoría y
compensación vinculada al ROIC a largo plazo --- no a las reservas trimestrales de ARR ---
crea la arquitectura de supervisión que mantiene unida la lógica de incentivos de Kerr. Alinear
los incentivos del consejo, la dirección y los colaboradores individuales con la métrica que
unifica todo el sistema --- el ROIC, o su proxy por cliente LTV:CAC --- resuelve la
desalineación de forma estructural en lugar de mediante supervisión gerencial.

%% ─────────────────────────────────────────────────────────────
\section{El fundador: influencia, narrativa y equipo}\label{sec:capstone-founder}
%% ─────────────────────────────────────────────────────────────

Cada número anterior es producido por personas a las que el fundador debe persuadir, financiar
y liderar. La Parte~VII trata esas capacidades no como habilidades blandas sino como insumos
para la misma prueba ROIC --- cada una mueve una variable que los capítulos de finanzas
mantienen fija.

La persuasión y la negociación (\cref{ch:persuasion-influence}) fijan dónde, dentro del rango
de negociación, se captura el valor. El término de la Serie~A es una negociación cuya zona de
posible acuerdo (ZOPA, \cref{eq:zopa}) va desde el valor de reserva del fundador
(\cref{eq:reservation}) --- el piso DCF de \money{15000000} por debajo del cual captar fondos
destruye más propiedad de la que compra --- hasta el techo del inversor cerca del caso
optimista de \money{28000000}. Dónde se fija el precio de la ronda dentro de esa banda es
excedente distributivo puro: un fundador que ancla en el comparable de \(6\times\) ARR en lugar
del caso pesimista retiene millones en propiedad que ninguna mejora operativa podría igualar.
La ecuación de confianza del mismo capítulo (\cref{eq:trust}) rige el lado de la demanda ---
los compradores que llegan pre-confiados a través de referencias convierten por debajo del
objetivo de CAC de \money{600} --- y es la razón por la que el género manipulador de las «leyes
del poder» se rechaza de plano: la confianza erosionada eleva el \emph{churn} y colapsa el
propio CLV sobre el que descansa todo el modelo.

La narrativa (\cref{ch:storytelling-communication}) es el mecanismo de entrega de esa persuasión.
El \emph{pitch} que convierte el piso de \money{15000000} en una pre-money financiada de
\money{24000000} es una narrativa, no una hoja de cálculo leída en voz alta; la misma disciplina
narrativa eleva la conversión de prueba a pago en la puerta vinculante del \emph{funnel} de
\cref{eq:funnel-chain} y siembra el boca a boca que reduce el CAC mixto. Una propuesta de valor
real que no puede transmitirse es, para el mercado, indistinguible de una que no existe.

La psicología del fundador (\cref{ch:founder-psychology}) es la disciplina que mantiene unido
el plan a lo largo de los ocho meses de \emph{runway}, y sus afirmaciones son las que sobreviven
al metaanálisis: la autoeficacia predice el rendimiento, los objetivos específicos y difíciles
superan a la intención vaga, el capital humano relevante para la tarea predice el éxito del
emprendimiento más que la experiencia genérica, y los hábitos de intención de implementación
convierten el propósito en rutina. Estos son precisamente los que sostienen la disciplina de
ejecución que \cref{eq:quick-ratio} y \cref{eq:burn-multiple} miden. El valor de empresa de
\money{12800000} se construye a partir de curvas de supervivencia de cohortes, densidad de señal
en la fase de aprendizaje y una distribución de escenarios explícita --- lo realiza un fundador
que instala esas rutinas, no uno que visualiza el resultado.

El liderazgo (\cref{ch:leadership-charisma}) cierra el bucle de vuelta al ROIC a través de la
plantilla. El liderazgo transformacional y las tácticas de carisma entrenables que
\cref{ch:storytelling-communication} suministra reclutan y retienen a los cinco ejecutivos de
cuenta que asume el modelo operativo de \cref{ch:operating-model}; la seguridad psicológica
determina si el hallazgo de reducción de \emph{churn} de la sección de ejecución se implementa
o se ignora silenciosamente; y las decisiones del dilema del fundador sobre los repartos de
\emph{equity} de cofundadores establecen la \emph{cap table} que \cref{eq:safe-conversion} diluye
después. Un equipo que no puede ser liderado con \money{4000000} de ARR no puede entregar el
plan que está valorando la pre-money de \money{24000000}. El liderazgo no es adyacente al
modelo financiero --- es el mecanismo que produce el NOPAT.

%% ─────────────────────────────────────────────────────────────
\section{Datos, IA y ESG: los añadidos de 2026}\label{sec:capstone-new}
%% ─────────────────────────────────────────────────────────────

Los tres temas añadidos con mayor prominencia en la edición de 2026 --- los datos como activo
estratégico (\cref{sec:data-asset}), la gobernanza de IA (\cref{sec:ai-governance}) y el ESG
como disciplina de asignación de capital (\cref{sec:esg-capital}) --- no son agendas estratégicas
separadas. Se reducen a la misma prueba ROIC\,>\,WACC que el resto del libro aplica a cada
decisión. Un volante de datos es estratégicamente valioso solo si genera mejoras de producto que
mueven el NOPAT hacia arriba o permiten a la empresa defender una prima de precio que sostiene el
margen de contribución; si no hace ninguna de las dos, es un costo de infraestructura disfrazado
de \emph{moat}. La gobernanza de IA vale su sobrecarga si y solo si el valor esperado de los
errores detectados supera el costo de revisión --- precisamente el cálculo de NPV que el libro
aplica a cualquier decisión de inversión, y exactamente el cálculo que \cref{sec:ai-governance}
establece con el CLV como denominador del valor en riesgo. El gasto en ESG crea beneficio
económico solo a través de dos canales específicos: reducir el WACC bajando la prima de riesgo
que los inversores asignan a la disrupción regulatoria o de talento, o mejorar el NOPAT
reduciendo el \emph{churn} (retención de talento) o elevando el poder de fijación de precios
(confianza de marca). Con \money{4000000} de ARR, el canal dominante es la reducción del WACC
mediante la fricción en la captación de fondos de la Serie~A --- la calidad de los datos ESG es
una puerta de diligencia, no un ejercicio altruista. Los tres temas son estratégicos solo cuando
mueven el NOPAT, el WACC o el capital invertido. Cuando no lo hacen, son costos.

%% ─────────────────────────────────────────────────────────────
\section{La síntesis}\label{sec:capstone-synthesis}
%% ─────────────────────────────────────────────────────────────

La columna vertebral de creación de valor en \cref{fig:value-spine} no es una metáfora; es una
cadena causal. Se comienza por la izquierda.

\emph{Los Fundamentos} establecen la prueba: el ROIC debe superar el WACC. Esa prueba es
abstracta hasta que la empresa rellena sus propios números. Para esta empresa, el veredicto es
\money{270000} de beneficio económico anual --- positivo, pero no seguro. El diferencial es del
\qty{9}{\percent} hoy; la pregunta es si persiste.

\emph{La economía del cliente} da a la prueba su contenido operativo. El segmento de mercado
mediano se eligió porque su potencial de LTV (\money{3150}) proporciona suficiente margen sobre
un CAC de \money{600} para sostener una razón de 5,25× --- la expresión por cliente del ROIC\,>\,WACC.
El trabajo de marca reduce el CAC efectivo con el tiempo; la disciplina de precios mantiene la
estructura de margen que genera los \money{252}/año de contribución. Sin ambas, el CLV se
contrae y la razón cae hacia el piso de 3×.

\emph{La adquisición} convierte el modelo de economía de unidad en demanda real. El \emph{funnel}
entrega 200 nuevos clientes al mes al costo objetivo, pero solo porque la fase de aprendizaje de
la plataforma tiene suficiente densidad de señal para pujar de forma inteligente. La máquina de
adquisición es frágil: fragmentar el presupuesto entre diez conjuntos de anuncios, o fijar el
objetivo de CPA por debajo del mínimo de la fase de aprendizaje, y todo el motor se detiene. La
disciplina de atribución --- iROAS en lugar de ROAS de último clic --- garantiza que la eficiencia
reportada es real, no tomada prestada de la demanda orgánica.

\emph{La estrategia} determina cuánto duran los resultados económicos. Un producto sin
\emph{moat} gana su diferencial ROIC\,>\,WACC hasta que un competidor replica las características;
típicamente ese período es de dos a cuatro años. Los costos de cambio por integración y el
naciente efecto de red del mercado de proveedores amplían el período elevando el costo de la
deserción y capitalizando el valor con cada nuevo nodo. Estos no son extras convenientes; son el
mecanismo por el cual la tasa de crecimiento terminal del \qty{3}{\percent} en el DCF es
defendible en lugar de arbitraria. La estrategia fija $g$.

\emph{Las finanzas} descuentan todo el sistema a la tasa apropiada, traduciendo la historia
operativa en un valor de empresa. El DCF no es un modelo de valoración; es una disciplina que
fuerza cada supuesto sobre crecimiento, margen y durabilidad en un único número que puede
someterse a pruebas de estrés. El análisis de escenarios revela que el caso pesimista ---
la regulación colapsa el \emph{moat} de costos de cambio, el crecimiento se detiene, el EV cae
a \money{8000000} --- es un evento con probabilidad del \qty{25}{\percent}. Gestionar ese riesgo
explícitamente, en lugar de asumir el caso base, es la diferencia entre las finanzas como función
de reporte y las finanzas como herramienta estratégica.

\emph{La ejecución} es donde la estrategia y las finanzas se realizan o no. El \emph{quick ratio}
de 1,85 y el múltiplo de \emph{burn} de 1,25 no son calificaciones; son coordenadas. Ubican al
negocio en la frontera de eficiencia y revelan la siguiente restricción: la reducción del
\emph{churn}, no el volumen de nueva adquisición, es la palanca que mueve ambas métricas hacia
territorio élite. La Ley de Little dice que el ciclo de ventas de 4 meses puede reducirse a 2
meses, pero solo reduciendo el \emph{pipeline} no calificado, no añadiendo vendedores a un
\emph{pipeline} que el equipo de cierre no puede liquidar más rápido. Las estructuras de
incentivos deben reforzar estos hallazgos, o la organización optimizará la métrica por la que
se le paga --- que rara vez es la métrica que importa.

\emph{El fundador} es la variable que el modelo no puede mantener constante. La persuasión fija
dónde en el rango de negociación se fija el precio de la ronda; la narrativa decide si el
\emph{pitch} y el \emph{funnel} convierten; la psicología suministra la disciplina --- y el
rechazo del pensamiento ilusorio --- que mantiene el pronóstico honesto; el liderazgo recluta y
retiene al equipo que convierte el plan en NOPAT. Estos no son un epílogo más blando a los
capítulos financieros; son los insumos humanos de cada número que la columna vertebral sostiene.

El bucle completo se cierra en el ROIC. La posición determina si los clientes están dispuestos
a pagar el precio que genera el margen de contribución. La economía del cliente determina si el
costo de adquisición es recuperable dentro de la vida del producto. La eficiencia de adquisición
determina si el CAC es realmente \money{600} o silenciosamente \money{900}. La durabilidad
estratégica determina si el margen y el crecimiento persisten lo suficiente para justificar el
valor terminal que domina el DCF. La disciplina de ejecución determina si el plan llega a tiempo.
Cada instrumento del libro apunta a la misma lectura: \(\mathrm{ROIC}>\mathrm{WACC}\) ---
sostenido, no solo hoy, sino a través del horizonte terminal que el mercado está valorando.

\textcite{brealey2020}, \textcite{damodaran2012}, \textcite{gupta2004customer},
\textcite{porter1985} y \textcite{osterwalder2010} son los cinco pilares cuyos \emph{frameworks}
este capítulo entrelaza. La integración es la contribución del libro; las fuentes son de ellos.
